% preamble.tex - All packages and styling for Python for Scientists
% ============================================================================

% Document class: KOMA-Script book (professional European book layout)
\documentclass[
    11pt,
    a4paper,
    twoside,
    open=right,
    chapterprefix=true,
    headings=normal,
    bibliography=totoc,
    listof=totoc
]{scrbook}

% ============================================================================
% ENCODING AND FONTS
% ============================================================================
\usepackage[T1]{fontenc}
\usepackage[utf8]{inputenc}
\usepackage{lmodern}            % Latin Modern fonts
\usepackage[expansion=false]{microtype}  % Improved typography (no expansion)

% ============================================================================
% PAGE LAYOUT
% ============================================================================
\usepackage[
    inner=2.5cm,
    outer=2cm,
    top=2.5cm,
    bottom=2.5cm,
    headheight=14pt
]{geometry}

% ============================================================================
% COLORS
% ============================================================================
\usepackage{xcolor}

% Syntax highlighting colors (adapted from original crash course)
\definecolor{codegreen}{rgb}{0.0, 0.5, 0.0}
\definecolor{codepurple}{rgb}{0.58, 0.0, 0.82}
\definecolor{codeorange}{rgb}{0.8, 0.4, 0.0}
\definecolor{codeblue}{rgb}{0.0, 0.0, 0.7}
\definecolor{codegray}{rgb}{0.5, 0.5, 0.5}
\definecolor{codecyan}{rgb}{0.0, 0.5, 0.5}
\definecolor{backcolour}{rgb}{0.97, 0.97, 0.97}

% Admonition box colors
\definecolor{warningbg}{RGB}{255, 243, 205}
\definecolor{warningborder}{RGB}{255, 100, 7}
\definecolor{tipbg}{RGB}{209, 236, 241}
\definecolor{tipborder}{RGB}{23, 162, 184}
\definecolor{notebg}{RGB}{232, 245, 233}
\definecolor{noteborder}{RGB}{76, 175, 80}
\definecolor{osbg}{RGB}{232, 234, 246}
\definecolor{osborder}{RGB}{63, 81, 181}
\definecolor{labbg}{RGB}{255, 248, 225}
\definecolor{labborder}{RGB}{150, 100, 50}
\definecolor{ponderbg}{RGB}{243, 229, 245}
\definecolor{ponderborder}{RGB}{156, 39, 176}

% ============================================================================
% CODE LISTINGS
% ============================================================================
\usepackage{listings}

% Python style
\lstdefinestyle{pythonstyle}{
    language=Python,
    backgroundcolor=\color{backcolour},
    commentstyle=\color{codegreen}\itshape,
    keywordstyle=\color{codeblue}\bfseries,
    numberstyle=\tiny\color{codegray},
    stringstyle=\color{codeorange},
    emphstyle=\color{codepurple},
    basicstyle=\ttfamily\small,
    breakatwhitespace=false,
    breaklines=true,
    captionpos=b,
    keepspaces=true,
    numbers=left,
    numbersep=8pt,
    showspaces=false,
    showstringspaces=false,
    showtabs=false,
    tabsize=4,
    frame=single,
    framesep=3pt,
    rulecolor=\color{codegray},
    xleftmargin=1.5em,
    framexleftmargin=1.5em,
    morekeywords={True, False, None, self, as, with, yield, lambda, assert},
    emph={print, range, len, type, isinstance, int, float, str, bool, list, dict, set, tuple, open, input, sum, min, max, abs, round, sorted, enumerate, zip, map, filter},
    escapeinside={(*@}{@*)},
    literate=
        {>>>}{{{\color{codecyan}>>>}}}3
        {...}{{{\color{codecyan}...}}}3
}

% Shell/terminal style
\lstdefinestyle{shellstyle}{
    backgroundcolor=\color{backcolour},
    basicstyle=\ttfamily\small,
    breakatwhitespace=false,
    breaklines=true,
    keepspaces=true,
    showspaces=false,
    showstringspaces=false,
    showtabs=false,
    tabsize=4,
    frame=single,
    framesep=3pt,
    rulecolor=\color{codegray},
    xleftmargin=0.5em,
    framexleftmargin=0.5em,
}

% Output style (no line numbers, slightly different background)
\lstdefinestyle{outputstyle}{
    backgroundcolor=\color{white},
    basicstyle=\ttfamily\small,
    breaklines=true,
    frame=single,
    framesep=3pt,
    rulecolor=\color{codegray!50},
    xleftmargin=0.5em,
    framexleftmargin=0.5em,
}

% Set Python as default
\lstset{style=pythonstyle}

% Inline code command
\newcommand{\py}[1]{\lstinline[style=pythonstyle]|#1|}
\newcommand{\shell}[1]{\lstinline[style=shellstyle]|#1|}

% ============================================================================
% ADMONITION BOXES
% ============================================================================
\usepackage{tcolorbox}
\tcbuselibrary{skins}

% Warning box
\newtcolorbox{warning}[1][]{
    enhanced,
    colback=warningbg,
    colframe=warningborder,
    coltitle=black,
    fonttitle=\bfseries,
    title={\textbf{Warning}},
    left=8pt,
    right=8pt,
    top=6pt,
    bottom=6pt,
    boxrule=1.5pt,
    #1
}

% Tip box
\newtcolorbox{tip}[1][]{
    enhanced,
    colback=tipbg,
    colframe=tipborder,
    coltitle=black,
    fonttitle=\bfseries,
    title={\textbf{Tip}},
    left=8pt,
    right=8pt,
    top=6pt,
    bottom=6pt,
    boxrule=1.5pt,
    #1
}

% Note box
\newtcolorbox{note}[1][]{
    enhanced,
    colback=notebg,
    colframe=noteborder,
    coltitle=black,
    fonttitle=\bfseries,
    title={\textbf{Note}},
    left=8pt,
    right=8pt,
    top=6pt,
    bottom=6pt,
    boxrule=1.5pt,
    #1
}

% OS-specific box
\newtcolorbox{osbox}[1][]{
    enhanced,
    colback=osbg,
    colframe=osborder,
    coltitle=white,
    fonttitle=\bfseries,
    title={\textbf{OS Differences}},
    left=8pt,
    right=8pt,
    top=6pt,
    bottom=6pt,
    boxrule=1.5pt,
    #1
}

% Lab example box (for chemistry/science examples)
\newtcolorbox{labexample}[1][]{
    enhanced,
    colback=labbg,
    colframe=labborder,
    coltitle=white,
    fonttitle=\bfseries,
    title={\textbf{Lab Example}},
    left=8pt,
    right=8pt,
    top=6pt,
    bottom=6pt,
    boxrule=1.5pt,
    #1
}

% Things to ponder box
\newtcolorbox{ponder}[1][]{
    enhanced,
    colback=ponderbg,
    colframe=ponderborder,
    coltitle=white,
    fonttitle=\bfseries,
    title={\textbf{Things to Ponder}},
    left=8pt,
    right=8pt,
    top=6pt,
    bottom=6pt,
    boxrule=1.5pt,
    #1
}

% ============================================================================
% MATHEMATICS
% ============================================================================
\usepackage{amsmath}
\usepackage{amssymb}

% ============================================================================
% GRAPHICS AND TABLES
% ============================================================================
\usepackage{graphicx}
\usepackage{booktabs}           % Professional tables
\usepackage{array}
\usepackage{multicol}

% ============================================================================
% HYPERLINKS AND CROSS-REFERENCES
% ============================================================================
\usepackage{hyperref}
\hypersetup{
    colorlinks=true,
    linkcolor=codeblue,
    filecolor=codepurple,
    urlcolor=codecyan,
    citecolor=codegreen,
    pdftitle={Python and Computational Thinking: A Crash Course for Scientists},
    pdfauthor={Frithjof Herb},
    pdfsubject={Python Programming for Scientists},
    pdfkeywords={Python, programming, science, chemistry, data analysis},
    bookmarksnumbered=true,
    bookmarksopen=true,
}

% ============================================================================
% HEADERS AND FOOTERS
% ============================================================================
\usepackage{scrlayer-scrpage}
\pagestyle{scrheadings}
\clearpairofpagestyles
\ohead{\pagemark}
\ihead{\headmark}

% ============================================================================
% MISCELLANEOUS
% ============================================================================
\usepackage{enumitem}           % Better list formatting
\setlist{noitemsep}             % Reduce spacing in lists

% Chapter formatting
\addtokomafont{chapter}{\color{codeblue}}
\addtokomafont{section}{\color{codeblue!80!black}}
\addtokomafont{subsection}{\color{codeblue!60!black}}

% Prevent widows and orphans
\clubpenalty=10000
\widowpenalty=10000

% ============================================================================
% CUSTOM COMMANDS
% ============================================================================

% Keyboard key styling
\newcommand{\key}[1]{\texttt{\fbox{\small #1}}}

% File/folder path styling
\newcommand{\filepath}[1]{\texttt{#1}}

% Menu navigation
\newcommand{\menu}[1]{\textsf{#1}}

% Variable names in running text
\newcommand{\var}[1]{\texttt{#1}}
